% Part: first-order-logic
% Chapter: syntax-and-semantics
% Section: formation-sequences

\documentclass[../../../include/open-logic-section]{subfiles}

\begin{document}

\olfileid{fol}{syn}{fseq}

\olsection{Formules stap voor stap opbouwen}

Met een inductieve definitie kun je !!{formula}s netjes en kort
opbouwen. Met inductie kun je daarna eigenschappen van !!{formula}s
bewijzen. Dat is een elegante en efficiënte aanpak. Soms is het echter
handiger om de bouw van !!{formula}s van onder naar boven, stap voor
stap, te volgen. Daarvoor gebruiken we hier een
\emph{constructiereeks}: een rij van opeenvolgende bouwstappen.
%
Zo kunnen we termen en !!{formula}s invoeren zonder steeds naar hun
inductieve definities te verwijzen. We beginnen met een willekeurige rij
tekens uit een taal~$\Lang L$. Zo'n rij noemen we een symboolrij.

\begin{defn}[Rijen tekens (symboolrijen)]
\ollabel{defn:string}
Neem een eerste-ordetaal $\Lang L$. Een \emph{$\Lang L$-symboolrij} is
een eindige rij tekens uit~$\Lang L$. Als uit de context duidelijk is
dat het om de taal~$\Lang L$ gaat, noemen we een
$\Lang L$-symboolrij kortweg een \emph{symboolrij}.
\end{defn}

\begin{ex}
Neem een willekeurige eerste-ordetaal $\Lang L$. Iedere
$\Lang L$-!!{formula} is dan een $\Lang L$-symboolrij. Andersom geldt
dat niet: \[)(\Obj v_0\lif\lexists\] is wel een
$\Lang L$-symboolrij, maar geen $\Lang L$-!!{formula}.
\end{ex}

\begin{defn}[Een term stap voor stap opbouwen]
\ollabel{defn:fseq-trm}
De eindige rij $\Lang L$-symboolrijen $\tuple{t_0,\dotsc,t_n}$ heet een
\emph{constructiereeks} voor een term $t$. Daarvoor moet zij met die
term eindigen: $t \ident t_n$. Voor iedere plaats $i \leq n$ moet bovendien een van
twee dingen gelden. Of $t_i$ is !!a{variable} of !!a{constant}. Of
$\Lang L$ bevat een $k$-plaatsige !!{function}~$f$ en er
bestaan eerdere plaatsen $m_0,\dotsc,m_k < i$ waarvoor
$t_i \ident f(t_{m_0},\dotsc,t_{m_k})$. Als we onderscheid moeten
maken, noemen we zo'n rij een \emph{termconstructiereeks}.
\end{defn}

\begin{ex}
De rij
\[
    \tuple{\Obj c_0, \Obj v_0, \Atom{\Obj f^2_0}{\Obj c_0, \Obj v_0}, \Atom{\Obj f^1_0}{\Atom{\Obj f^2_0}{\Obj c_0, \Obj v_0}}}
\]
bouwt stap voor stap de term $\Atom{\Obj f^1_0}{\Atom{\Obj
f^2_0}{\Obj c_0, \Obj v_0}}$ op en is dus een constructiereeks voor
die term. Deze rij doet dat ook:
\[
    \tuple{\Obj v_0, \Obj c_0, \Atom{\Obj f^2_0}{\Obj c_0, \Obj v_0}, \Atom{\Obj f^1_0}{\Atom{\Obj f^2_0}{\Obj c_0, \Obj v_0}}}.
\]
\end{ex}

\begin{defn}[Een formule stap voor stap opbouwen]
\ollabel{defn:fseq-frm}
De eindige rij $\Lang L$-symboolrijen $\tuple{!A_0,\dotsc,!A_n}$
heet een \emph{constructiereeks} voor~$!A$. Haar laatste element moet
de bedoelde formule zijn: $!A \ident !A_n$. Voor iedere plaats
$i \leq n$ geldt één van twee dingen. Het element $!A_i$ is een
atomaire !!{formula}. Of het is opgebouwd uit eerdere elementen.
Preciezer: er bestaan dan $j,k < i$ en
!!a{variable}~$x$ waarvoor een van deze regels geldt:
\begin{enumerate}
    \tagitem{prvNot}{$!A_i \ident \lnot !A_j$.}{}%
    \tagitem{prvAnd}{$!A_i \ident (!A_j \land !A_k)$.}{}%
    \tagitem{prvOr}{$!A_i \ident (!A_j \lor !A_k)$.}{}%
    \tagitem{prvIf}{$!A_i \ident (!A_j \lif !A_k)$.}{}%
    \tagitem{prvIff}{$!A_i \ident (!A_j \liff !A_k)$.}{}%
    \tagitem{prvAll}{$!A_i \ident \lforall[x][!A_j]$.}{}%
    \tagitem{prvEx}{$!A_i \ident \lexists[x][!A_j]$.}{}%
\end{enumerate}
Als we onderscheid moeten maken, noemen we zo'n rij een
\emph{formuleconstructiereeks}.
\end{defn}

\begin{ex}
\[
\tuple{
    \Atom{\Obj A^1_0}{\Obj v_0},
    \Atom{\Obj A^1_1}{\Obj c_1},
    (\Atom{\Obj A^1_1}{\Obj c_1} \land \Atom{\Obj A^1_0}{\Obj v_0}),
    \lexists[\Obj v_0][(\Atom{\Obj A^1_1}{\Obj c_1} \land \Atom{\Obj A^1_0}{\Obj v_0})]
}
\]
is een constructiereeks voor $\lexists[\Obj v_0][(\Atom{\Obj A^1_1}{\Obj
c_1} \land \Atom{\Obj A^1_0}{\Obj v_0})]$. De volgende rij is dat
ook:
\begin{multline*}
\tuple{
    \Atom{\Obj A^1_0}{\Obj v_0},
    \Atom{\Obj A^1_1}{\Obj c_1},
    (\Atom{\Obj A^1_1}{\Obj c_1} \land \Atom{\Obj A^1_0}{\Obj v_0}),
    \Atom{\Obj A^1_1}{\Obj c_1},\\
    \lforall[\Obj v_1][\Atom{\Obj A^1_0}{\Obj v_0}],
    \lexists[\Obj v_0][(\Atom{\Obj A^1_1}{\Obj c_1} \land \Atom{\Obj A^1_0}{\Obj v_0})]
}.
\end{multline*}
%
In de tweede rij staan overbodige stappen. Sommige !!{formula}s zijn
dubbel of zijn niet nodig om de eindformule te bouwen. Zulke stappen
zijn de ``rommel'' die een constructiereeks mag bevatten.
\end{ex}

\begin{prop}\ollabel{prop:formed}
Iedere !!{formula}~$!A$ in~$\Frm[L]$ heeft een constructiereeks.
\end{prop}

\begin{proof}
Is $!A$ atomair, dan is de rij $\tuple{!A}$ al een
constructiereeks voor~$!A$.
%
Neem nu aan dat $!B$ en~$!C$ constructiereeksen hebben. Noem die
respectievelijk $\tuple{!B_0,\dotsc,!B_n}$ en
$\tuple{!C_0,\dotsc,!C_m}$.
%
\begin{enumerate}
  \tagitem{prvNot}{Als $!A \ident \lnot !B$, voeg dan de negatie van
      het laatste element achter de bijbehorende rij. De rij
      $\tuple{!B_0,\dotsc,!B_n,\lnot !B_n}$
      is een constructiereeks voor~$!A$.}{}
  \tagitem{prvAnd}{Als $!A \ident (!B \land !C)$, zet dan de twee
      rijen achter elkaar en voeg de laatste bouwstap toe. De rij
      $\tuple{!B_0,\dotsc,!B_n,!C_0,\dotsc,!C_m,(!B_n \land !C_m)}$
      is een constructiereeks voor~$!A$.}{}
  \tagitem{prvOr}{Neem $!A \ident (!B \lor !C)$. Dan is
      $\tuple{!B_0,\dotsc,!B_n,!C_0,\dotsc,!C_m,(!B_n \lor !C_m)}$
      een constructiereeks voor~$!A$.}{}
  \tagitem{prvIf}{Neem $!A \ident (!B \lif !C)$. Dan is
      $\tuple{!B_0,\dotsc,!B_n,!C_0,\dotsc,!C_m,(!B_n \lif !C_m)}$
      een constructiereeks voor~$!A$.}{}
  \tagitem{prvIff}{Neem $!A \ident (!B \liff !C)$. Dan is
      $\tuple{!B_0,\dotsc,!B_n,!C_0,\dotsc,!C_m,(!B_n \liff !C_m)}$
      een constructiereeks voor~$!A$.}{}
  \tagitem{prvAll}{Neem $!A \ident \lforall[x][!B]$. Dan is
      $\tuple{!B_0,\dotsc,!B_n,\lforall[x][!B_n]}$
      een constructiereeks voor~$!A$.}{}
  \tagitem{prvEx}{Neem $!A \ident \lexists[x][!B]$. Dan is
      $\tuple{!B_0,\dotsc,!B_n,\lexists[x][!B_n]}$
      een constructiereeks voor~$!A$.}{}
\end{enumerate}
Dit zijn alle manieren waarop !!{formula}s worden opgebouwd. Volgens
het inductieprincipe heeft daarom iedere !!{formula} een
constructiereeks.
\end{proof}

Ook de omgekeerde richting geldt: iedere symboolrij met een
formuleconstructiereeks is werkelijk een formule. Onze twee manieren
om formules te definiëren leveren dus precies dezelfde uitkomsten op.
Daardoor kunnen we stellingen over formules ook bewijzen met gewone
inductie naar de lengte van een constructiereeks.

\begin{lem}
\ollabel{lem:fseq-init}
Neem een constructiereeks $\tuple{!A_0,\dotsc,!A_n}$ voor~$!A_n$ en
neem $k \leq n$. Stoppen we de rij op die plaats, dan is de kortere rij
$\tuple{!A_0,\dotsc,!A_k}$ een constructiereeks voor~$!A_k$.
\end{lem}

\begin{proof}
Opgave.
\end{proof}

\begin{prob}
Bewijs \olref[fol][syn][fseq]{lem:fseq-init}.
\end{prob}

\begin{thm}
\ollabel{thm:fseq-frm-equiv}
$\Frm[L]$ bestaat precies uit de $\Lang L$-symboolrijen~$!A$
waarvoor een formuleconstructiereeks voor~$!A$ bestaat.
\end{thm}

\begin{proof}
Noem $F$ de verzameling van alle symboolrijen uit de taal~$\Lang L$ met
een constructiereeks. Uit
\olref[fol][syn][fseq]{prop:formed} weten we al dat
$\Frm[L] \subseteq F$: iedere formule behoort dus tot die verzameling. Nu bewijzen we
de andere richting.

Neem een symboolrij $!A$ met een constructiereeks
$\tuple{!A_0,\dotsc,!A_n}$. We bewijzen dat $!A \in \Frm[L]$ met
sterke inductie naar~$n$. De inductiehypothese zegt dat iedere symboolrij
met een constructiereeks van lengte $m < n$ al in $\Frm[L]$ zit.
Volgens de definitie van een constructiereeks is $!A \ident !A_n$
atomair, of er zijn eerdere plaatsen $j,k < n$ waarvoor een van deze
gevallen geldt:
\begin{enumerate}
    \tagitem{prvNot}{$!A \ident \lnot !A_j$.}{}%
    \tagitem{prvAnd}{$!A \ident (!A_j \land !A_k)$.}{}%
    \tagitem{prvOr}{$!A \ident (!A_j \lor !A_k)$.}{}%
    \tagitem{prvIf}{$!A \ident (!A_j \lif !A_k)$.}{}%
    \tagitem{prvIff}{$!A \ident (!A_j \liff !A_k)$.}{}%
    \tagitem{prvAll}{$!A \ident \lforall[x][!A_j]$.}{}%
    \tagitem{prvEx}{$!A \ident \lexists[x][!A_j]$.}{}%
\end{enumerate}
We bekijken de gevallen apart. Is $!A$ atomair, dan is
$!A_n \in \Frm[L_0]$. Neem nu het geval
$!A \equiv (!A_j \land !A_k)$. Volgens
\olref[fol][syn][fseq]{lem:fseq-init} zijn
$\tuple{!A_0,\dotsc,!A_j}$ en $\tuple{!A_0,\dotsc,!A_k}$
constructiereeksen. De eerste is voor $!A_j$ en de tweede voor~$!A_k$.
Beide kortere rijen beginnen bij hetzelfde eerste element en stoppen
eerder dan de rij voor~$!A$. Hun lengte is dus kleiner dan~$n$.
Volgens de inductiehypothese zitten $!A_j$ en~$!A_k$
daarom in~$\Frm[L_0]$. Uit de definitie van !!a{formula} volgt dan dat
ook $(!A_j \land !A_k)$ daarin zit. De andere gevallen gaan op dezelfde
manier.
\end{proof}

Voor constructiereeksen van termen gelden vergelijkbare eigenschappen
als voor die van !!{formula}s.

\begin{prop}
\ollabel{prop:fseq-trm-equiv}
$\Trm[L]$ bestaat precies uit de $\Lang L$-symboolrijen $t$
waarvoor een termconstructiereeks voor~$t$ bestaat.
\end{prop}

\begin{proof}
Opgave.
\end{proof}

\begin{prob}
Bewijs \olref[fol][syn][fseq]{prop:fseq-trm-equiv}.
Aanwijzing: gebruik ongeveer dezelfde aanpak als in het bewijs van
\olref[fol][syn][fseq]{thm:fseq-frm-equiv}.
\end{prob}

Een constructiereeks kan op twee manieren overbodige stappen bevatten.
Soms komt hetzelfde element meer dan één keer voor. Of een element kan nergens nodig zijn
om de uiteindelijke formule of term op te bouwen. Beide soorten
``rommel'' verdwijnen als we alleen naar de kortst mogelijke
constructiereeksen kijken.

\begin{defn}[Kortst mogelijke constructiereeksen]
\ollabel{defn:minimal-fseq}
Een constructiereeks $\tuple{!A_0, \ldots, !A_n}$ voor een
formule~$!A$ is een \emph{kortst mogelijke constructiereeks} voor~$!A$
als geen andere constructiereeks met hetzelfde einddoel korter is.
Neem formeel een willekeurige andere constructiereeks~$s$ voor~$!A$.
De lengte van~$s$ is dan groter dan of gelijk aan~$n+1$.

Zo is ook een constructiereeks $\tuple{t_0, \ldots, t_n}$
voor een term~$t$ een \emph{kortst mogelijke constructiereeks}
voor~$t$ als geen andere constructiereeks met hetzelfde einddoel korter
is. Neem een willekeurige andere constructiereeks~$s$ voor~$t$. De
lengte van~$s$ is dan groter dan of gelijk aan~$n+1$.
\end{defn}

Kortst mogelijk betekent niet dat er maar één zo'n rij is. Een formule
of term kan meer dan één kortst mogelijke constructiereeks hebben. Al
die reeksen bevatten wel precies dezelfde symboolrijen.

\begin{prop}
\ollabel{prop:subformula-equivs}
De volgende drie uitspraken zijn equivalent; ze zeggen dus hetzelfde:
\begin{enumerate}
    \item $!B$ is een sub-!!{formula} van~$!A$.
    \item $!B$ komt voor in iedere constructiereeks van~$!A$.
    \item $!B$ komt voor in een kortst mogelijke constructiereeks van~$!A$.
\end{enumerate}
\end{prop}

\begin{proof}
Opgave.
\end{proof}

\begin{prob}
Bewijs \olref[fol][syn][fseq]{prop:subformula-equivs}.
\end{prob}

\begin{history}
Raymond Smullyan voerde constructiereeksen in in zijn leerboek
\emph{First-Order Logic} \citep{Smullyan1968}.
\citet{Zuckerman1973} stelde later meer eigenschappen van
constructiereeksen vast.
\end{history}
\end{document}
